\documentclass[11pt,a4paper]{article}
\usepackage[T1]{fontenc}
\usepackage[utf8]{inputenc}
\usepackage{amsmath,amssymb,amsthm}
\usepackage[margin=1in]{geometry}
\usepackage[colorlinks=true,linkcolor=blue,urlcolor=blue]{hyperref}
\theoremstyle{plain}
\newtheorem{theorem}{Theorem}
\newtheorem{lemma}[theorem]{Lemma}
\newtheorem{proposition}[theorem]{Proposition}
\newtheorem{corollary}[theorem]{Corollary}
\theoremstyle{definition}
\newtheorem{remark}[theorem]{Remark}

\title{The empirical recurrence for OEIS A204498, proved by transfer matrix}
\author{Adrian Perez Fontelles\\ \small Independent researcher}
\date{7 September 2026}
\begin{document}
\maketitle

\begin{abstract}
OEIS A204498 counts $(n+2)\times7$ arrays over $\{0,\dots,2\}$ subject to a condition imposed on
every $3\times3$ subblock, and it carries an empirical recurrence of order $64$
contributed by R.~H.~Hardin. It is true, and it is not an empirical matter. A $3\times3$
subblock spans three consecutive rows, so the state of a transfer matrix is a PAIR of
consecutive rows and appending one more row is a step; an array is then a walk, and
$a(n)$ is a walk count on $S=57776$ vertices. Such a sequence is $C$-finite, so whether a given
recurrence annihilates it is decided by a finite exact computation, with the Cayley--Hamilton
theorem bounding the work.
\end{abstract}

\noindent\small 2020 Mathematics Subject Classification. 05A15, 05B45, 15B36.\normalsize

\section{The conjecture}

OEIS A204498 is ``Number of (n+2) X 7 0..2 arrays with every 3 X 3 subblock having three strictly increasing elements in a row horizontally or vertically.''. It has offset $1$ and begins
\[
65606908, 493398441, 3151347159, 22441401117, 208673819127, 2160920648295, 23772717607473, 276423235451370,\ \dots
\]
The entry states, as a conjecture contributed by its author and never marked settled:
\begin{quote}
Empirical: a(n) = 1459*a(n-3) +1966*a(n-4) +4770*a(n-5) -506169*a(n-6) -729427*a(n-7) -3282546*a(n-8) -3675924*a(n-9) +593623566*a(n-10) +951185496*a(n-11) -448513794*a(n-12) -1134670248*a(n-13) -109937354733*a(n-14) -77114079432*a(n-15) +279521405685*a(n-16) -459219944592*a(n-17) +4495749207192*a(n-18) -104133790895319*a(n-19) -323002370068452*a(n-20) -904277651409864*a(n-21) -1834603331470920*a(n-22) +8183503227524310*a(n-23) +53104054207803750*a(n-24) +160385257720234269*a(n-25) +27116289266862438*a(n-26) -358008897750745914*a(n-27) -2323633608815164830*a(n-28) -38965319111379563784*a(n-29) -45315387799040184984*a(n-30) -71831628696486933783*a(n-31) +961949802637047983886*a(n-32) +1928931499438085567358*a(n-33) +11098386323688924911358*a(n-34) -2781662912179647932907*a(n-35) -18944555420780928351*a(n-36) -180443740883973269163453*a(n-37) -174757345526978402459058*a(n-38) -1142655081431618149373196*a(n-39) -283873257921844853861508*a(n-40) -1758486400193961357313203*a(n-41) +5418976385678390064525933*a(n-42) -4333456233936328096284678*a(n-43) +52357850708967075345352317*a(n-44) +113875738890358473567768429*a(n-45) +332295902293018960798375482*a(n-46) +507808614055252846221859563*a(n-47) +1039268121466197777707272662*a(n-48) -1971166572615872347630335675*a(n-49) -4385193296346556587211012596*a(n-50) +3197026868983413864936749091*a(n-51) +4167793366663358828991037653*a(n-52) +1890644766282248394058214985*a(n-53) +7508799720706674292816020399*a(n-54) +2233085916306206014836593265*a(n-55) -3034648478031704823099463659*a(n-56) -1006738221464962955098662348*a(n-57) -459404171909369698783709196*a(n-58) +2414415236064884629304721732*a(n-59) +2254779373005991286687286144*a(n-60) +190196166013621067622528648*a(n-61) +64499687123220301474143972*a(n-62) +17220705727178598190290720*a(n-63) -175070382610419774734400*a(n-64) for n>75.
\end{quote}
As of the ``Last modified'' line on the live entry (Sep 05 2025 13:41:02, revision 9) this is still
recorded as empirical, and nothing on the entry records it as proved.

\section{Three consecutive lines}

Write an admissible array as a sequence of rows $u_1,u_2,\dots$, each an element of
$\{0,\dots,2\}^{7}$. A $3\times3$ subblock occupies three consecutive rows and
three consecutive columns, so with $r,s,t$ three consecutive rows the blocks it
contributes are
\[
g=\begin{pmatrix}r_j&r_{j+1}&r_{j+2}\\ s_j&s_{j+1}&s_{j+2}\\ t_j&t_{j+1}&t_{j+2}\end{pmatrix}
\]
for each admissible index, and the entry's condition is the requirement that
\[
\text{at least one}\text{ of the horizontally or vertically lines of }g\text{ \text{is strictly increasing}}
\]
holds for every such $g$. This involves $r$, $s$ and $t$ only: it is a condition on
three consecutive rows and on nothing else.

Take as vertices of a digraph $G$ the ordered PAIRS $(r,s)$ of rows, of which there are
$(2+1)^{2\cdot7}=4782969$, and put an edge $(r,s)\to(s,t)$ exactly when the condition
holds for the triple $(r,s,t)$ at every admissible index. An array with $L$ rows is then a
walk of length $L-2$ in $G$: its first two rows name the starting vertex and each further
row is one step. Most of those pairs are indistinguishable, and exactly so. For a window $j$, whether a row $t$ is admissible there depends on $r$ and $s$ only through their own $j$-th windows, so collect for each $j$ the set of triples $(t_j,t_{j+1},t_{j+2})$ the pair allows and call that tuple $C(r,s)$. The rows that may follow $(r,s)$ are exactly those consistent with every $C_j$, and the successor $(s,t)$ carries $C(s,t)$, which does not mention $r$. Hence two pairs $(r,s)$ and $(r',s)$ with the same $C$ have the same outgoing rows and the same successors: they may be identified, and the vertex is $(s,C)$. Each such vertex begins an array with the multiplicity of $r$ that produced it, which is what $\iota$ records. That leaves $S=57776$ vertices. Hence, with $M$ the adjacency matrix of $G$ on $S=57776$ vertices and
$\iota$ the vector of starting weights,
\[
a(n)\;=\;\,\iota^{\!\top}M^{\,n-1+1}\mathbf{1},
\]
the exponent being fixed by the entry's own shape and checked against its published terms.
In particular $a$ is $C$-finite.

\section{The proof}

\begin{lemma}\label{lem:crit}
Let $q(t)=t^{64}-\sum_i c_it^{64-i}$ be the characteristic polynomial of the
conjectured recurrence. Then the residual $a(n)-\sum_ic_ia(n-i)$ equals
$\mathbf{1}^{\!\top}M^{\,j}q(M)\mathbf{1}$ for the corresponding walk index
$j$, so the recurrence holds from some point on if and only if
$\mathbf{1}^{\!\top}M^{j}q(M)\mathbf{1}=0$ for all large $j$; and by the
Cayley--Hamilton theorem it is enough to examine $j<S$.
\end{lemma}

\begin{proof}
Factor $M^{j}$ out of $M^{j+64}-\sum_ic_iM^{j+64-i}$. For the bound, $M^{S}$
is an integer combination of $I,M,\dots,M^{S-1}$, so the numbers
$u_j=\mathbf{1}^{\!\top}M^{j}q(M)\mathbf{1}$ satisfy the monic linear recurrence given by
the characteristic polynomial of $M$; $S$ consecutive zeros therefore force every later one,
and the last nonzero $u_j$ pins the threshold exactly.
\end{proof}

\begin{theorem}
\[
a(n)\;=\;1459\,a(n-3) + 1966\,a(n-4) + 4770\,a(n-5) - 506169\,a(n-6) - 729427\,a(n-7) - 3282546\,a(n-8) - 3675924\,a(n-9) + 593623566\,a(n-10) + 951185496\,a(n-11) - 448513794\,a(n-12) - 1134670248\,a(n-13) - 109937354733\,a(n-14) - 77114079432\,a(n-15) + 279521405685\,a(n-16) - 459219944592\,a(n-17) + 4495749207192\,a(n-18) - 104133790895319\,a(n-19) - 323002370068452\,a(n-20) - 904277651409864\,a(n-21) - 1834603331470920\,a(n-22) + 8183503227524310\,a(n-23) + 53104054207803750\,a(n-24) + 160385257720234269\,a(n-25) + 27116289266862438\,a(n-26) - 358008897750745914\,a(n-27) - 2323633608815164830\,a(n-28) - 38965319111379563784\,a(n-29) - 45315387799040184984\,a(n-30) - 71831628696486933783\,a(n-31) + 961949802637047983886\,a(n-32) + 1928931499438085567358\,a(n-33) + 11098386323688924911358\,a(n-34) - 2781662912179647932907\,a(n-35) - 18944555420780928351\,a(n-36) - 180443740883973269163453\,a(n-37) - 174757345526978402459058\,a(n-38) - 1142655081431618149373196\,a(n-39) - 283873257921844853861508\,a(n-40) - 1758486400193961357313203\,a(n-41) + 5418976385678390064525933\,a(n-42) - 4333456233936328096284678\,a(n-43) + 52357850708967075345352317\,a(n-44) + 113875738890358473567768429\,a(n-45) + 332295902293018960798375482\,a(n-46) + 507808614055252846221859563\,a(n-47) + 1039268121466197777707272662\,a(n-48) - 1971166572615872347630335675\,a(n-49) - 4385193296346556587211012596\,a(n-50) + 3197026868983413864936749091\,a(n-51) + 4167793366663358828991037653\,a(n-52) + 1890644766282248394058214985\,a(n-53) + 7508799720706674292816020399\,a(n-54) + 2233085916306206014836593265\,a(n-55) - 3034648478031704823099463659\,a(n-56) - 1006738221464962955098662348\,a(n-57) - 459404171909369698783709196\,a(n-58) + 2414415236064884629304721732\,a(n-59) + 2254779373005991286687286144\,a(n-60) + 190196166013621067622528648\,a(n-61) + 64499687123220301474143972\,a(n-62) + 17220705727178598190290720\,a(n-63) - 175070382610419774734400\,a(n-64)
\]
for every $n>75$.
\end{theorem}

\begin{proof}
The vector $w=q(M)\mathbf{1}$ was formed by $64$ matrix--vector products in exact
integer arithmetic and $\mathbf{1}^{\!\top}M^{j}w$ evaluated for $j=0,1,\dots$, again
exactly. Those integers vanish from the index corresponding to $n=75+1$ onwards, and the
run was continued until the working vector was identically zero, which settles all larger $j$
at once. By Lemma~\ref{lem:crit} the recurrence holds for every $n>75$.
\end{proof}

\section{Verification}

Three checks, each independent of the algebra above.

First, the transfer model reproduces all $13$ terms the entry publishes, exactly, in
integer arithmetic. This is what ties the model to the entry: the digraph is built by reading
the entry's English, and a misreading gives different counts.

Second, the conjectured recurrence was evaluated directly on the entry's own published terms,
with no matrices involved, and holds wherever the proved range and the data overlap.

Third, the annihilation test of Lemma~\ref{lem:crit} was carried out in exact integer
arithmetic on the whole vector rather than on a sampled prefix, so no rounding or sampling
enters anywhere.

\begin{thebibliography}{9}
\bibitem{oeis} The OEIS Foundation, \emph{The On-Line Encyclopedia of Integer
Sequences}, \url{https://oeis.org}, 2026. Sequence A204498.
\bibitem{stanley} R.~P.~Stanley, \emph{Enumerative Combinatorics}, Volume 1, 2nd ed.,
Cambridge University Press, 2011. (Transfer-matrix method, Section 4.7.)
\bibitem{fs} P.~Flajolet and R.~Sedgewick, \emph{Analytic Combinatorics}, Cambridge
University Press, 2009.
\bibitem{horn} R.~A.~Horn and C.~R.~Johnson, \emph{Matrix Analysis}, 2nd ed., Cambridge
University Press, 2013.
\end{thebibliography}

\end{document}
